\section{Forensic Sequencing Details}
\label{app:forensic_sequence_v18}

The sequential architecture in the main text separates two tasks
that are often conflated. The award-layer stage ranks firms using
fields that are available before bid recovery. The bid-layer
stage then uses richer forensic features only for the firms that
survive the first stage. This appendix states the sequence
explicitly.

\subsection{Stage 1: Award-Layer Ranking}
\label{app:stage1_v18}

Stage 1 uses the always-loser stratum and ranks firms by
$\log(1+T_i)$. An enforcement agency can translate the rank into
a fixed threshold, a top-$k$ list, or a budget-constrained request
for further evidence. No bid values are needed at this stage.
The output is a survivor pool, not a liability finding.

\subsection{Stage 2: Bid-Layer Forensics}
\label{app:stage2_v18}

Stage 2 is run only inside the survivor pool. The bid-layer
benchmark uses within-tender bid-distribution features:
dispersion, skewness, kurtosis, spread, distance from the winning
bid, and related moments. These features are closer to forensic
proof because they describe how bids are placed inside tenders.
They are also more costly because they require per-bid offer
values and sufficient tender-level bid histories.

The comparison in the main text uses the bid-layer benchmark in
two roles. First, it is an external benchmark for the award-layer
screen: the cheap screen is informative because it reaches
similar discrimination on a thinner data envelope. Second, it is
the Stage 2 input in the sequential rule: the screen does not
replace bid-distribution analysis; it decides where that analysis
should begin.

\subsection{Why Joint Scoring Is Not the Operational Baseline}
\label{app:joint_scoring_v18}

A single joint model that uses both award-layer and bid-layer
features for every firm is a useful upper-bound comparison, but
it is not the relevant operational baseline for many agencies.
Running the joint model requires bid microdata for the full
evaluation pool before any triage has occurred. The sequential
rule asks a different question: how much of the joint model's
target capture can be preserved while recovering bid microdata
for only a smaller first-stage pool?

That distinction is the enforcement contribution. Joint scoring
is a prediction exercise under full observability. Sequencing is
an evidence-allocation rule under partial observability. The
second problem is the one agencies face when award records are
routine but bid files require costly recovery.

\subsection{Algorithm}
\label{app:gatekeeping_algorithm_v18}

The full gatekeeping procedure can be stated compactly.

\begin{quote}
\textbf{Algorithm F.1: Award-Layer Gatekeeping}
\begin{enumerate}
\item Define the firm universe from award records (participant
identity, winner identity, item code, negotiated price,
procuring unit, year, modality).
\item For each firm $i$, compute the win count $W_i$ and the
participation count $T_i$ over the calibration window.
\item Restrict to the zero-win stratum $\{i : W_i = 0\}$.
\item Compute the award-layer score $s_i = \log(1+T_i)$.
\item Select the survivor pool $\mathcal{S}$ using one of: a
median-plus-$1.5\times$IQR threshold, a top-$K_1$ rule, or a
capacity-constrained budget rule.
\item Recover bid microdata only for firms in $\mathcal{S}$.
\item Compute the bid-layer forensic features inside
$\mathcal{S}$.
\item Rank survivor firms using the bid-layer benchmark alone
or a combined award-plus-bid rule, depending on the
investigative protocol.
\item Deliver the top-$k$ queue from that ranking for
investigative review.
\item Treat the output as forensic priority, not as a liability
finding.
\end{enumerate}
\end{quote}

Joint scoring requires bid microdata for the full pool before
triage; sequential gatekeeping requires bid microdata only
after the award-layer survivor pool is formed. The size of the
survivor pool, the choice of threshold or top-$K_1$, and the
recalibration cadence are all institutional choices that the
algorithm leaves to the agency.
